\section{A better score can still lose the test}
\label{sec:self-forecast}

Keep the last measured value. That simple prediction can be hard to beat when nearby observations change slowly. A separate computer-learning diagnostic compared it with a model trained to forecast later muscle and glove features. More optimization improved the model's full-TRAIN score over its own earlier states. Carrying the last value forward still won in both streams. This is a different experiment from LINK and from the moving simulated hand.\cite{movementTrain990Recorded2026}

\subsection{What the model actually receives}
These observations come from Reach\&Grasp version 1.0 in IIT Dataverse. Its dataset metadata declares the Creative Commons Attribution 4.0 licence, or CC BY 4.0.\cite{reachGraspV1} The studied sub-03/Cyl pair consists of a cyberglove-motion file (provider ID 14486, 999,735 bytes) and a sessantaquattro-EMG file (ID 14393, 77,809,601 bytes). Their retained acquisition record reports checking each provider MD5, a file fingerprint. That is a recorded check, not a new reading or checksum performed for this paper revision. Identifying the files does not calibrate the sensors, establish their clock alignment or settle every participant and combined-use right. It does not turn one acquisition into twelve newly examined trajectories.

There is one original TRAIN acquisition, called sub-03/Cyl. Four invalid rows were excluded from 3,127 source rows, leaving 3,123. Each retained row contains 64 derived surface-EMG RMS features and 18 glove-angle coordinates. RMS means root mean square: square the signal samples in the specified trailing 40 millisecond interval, average those squared values and take the square root. The resulting feature summarizes that interval rather than a single instantaneous sample. The source declares millivolt-derived RMS and degree angles. Those declarations are not an independent calibration of the apparatus, its clocks or when a feature could actually become available. A trailing feature depends on more than the instant printed beside it.

Each stream has its own adapter, which puts its observations into the form the model receives, and its own output head, which predicts that stream's target features. Between them sits one shared gated recurrent unit (GRU), a sequence model whose changing numerical state summarizes earlier inputs. The muscle task sees 64 rows of muscle features and forecasts later muscle features. The glove task sees 64 rows of glove features and forecasts later glove features. Targets are one or ten rows later. Shared fitted numbers do not turn these into a muscle-to-motion or neural-to-body experiment. There is no recorded neural input, force target or sensory-afferent target in these tasks.

One and ten are positions in recorded row order. Without the original timing and availability qualification, they cannot be called physiological conduction delays or proof that a live predictor can respond in time. Surface EMG records muscle electrical activity; it does not directly record the incoming sensory nerves or every consequence of contraction.

\subsection{What changed during continuation}
The first pilot made 40 joint updates. An update is one change to the fitted numbers, not an epoch through the whole catalog. Each update used four ordered windows from each task, so only the first 160 windows per task supplied that pilot's optimization examples. On the full TRAIN comparison, its score was worse than initialization and the last-value predictor.

The continuation began from that saved 40-update state. It restored the model's weights and AdamW's saved running summaries and step 40, then restored the saved Torch random-generator state after constructing the temporary objects and loading the optimizer. Restoring only weights could produce a similar prediction while changing the next update; the saved update state matters to what was continued.

There were 950 further updates. The first 949 used 32 windows per task; the last used 17. Together with the earlier 160 exposures, this gives 30,545 window exposures per task across 990 cumulative updates: five passes through the 6,109-window catalog. Of those windows, 3,059 use the one-row target offset and 3,050 the ten-row offset. A window can overlap another window, and a later pass can repeat the same window. None becomes another person or another independent acquisition because it was used again.

\subsection{Give every condition the same question}
The comparison kept four conditions: the untrained initial state, the frozen 40-update state, the continued 990-update state and persistence. Persistence predicts each target feature by carrying forward its last supported past value. Every condition faced the same windows, targets, missing-value masks and saved original-TRAIN scalers. A scaler puts the observations on the specified numerical scale. Its fitted values came once from the original TRAIN frames; no new scaler was fitted in continuation or evaluation. Evaluation changed no model weights.

The score squares each prediction's difference from its normalized target, sums those squared errors and divides by the number of observed target features. If $z$ is the normalized target, $\widehat z$ its prediction and $m$ marks an observed target, the same rule is
\[
 \operatorname{MSE}_{\rm norm}=\frac{\sum m(\widehat z-z)^2}{\sum m}.
\]
The recorded predictions and residual coordinates use float32 numbers. Their squared errors are accumulated in float64, a format with more numerical precision, and divided by the exact count. Combining offsets adds their errors and their counts before division. The pooled counts are 390,976 RMS target features and 109,962 glove target features. These scores have no physical unit after normalization: they are neither squared millivolts nor squared degrees. More decimal places preserve the numerical record; they do not prove that the measurements or biological predictions are that accurate.

The original first 4 windows were also checked, two at each offset. They sit inside the full catalog and the initial optimization prefix. They are a small overlapping view of the same acquisition. They supply no new independent sample, population confidence interval or statistical $p$-value.

\input{../../common/training-result-tables.tex}

\subsection{Keep the failure beside the improvement}
The final RMS score is 0.48371488886038944, below the initial 1.0195127528126222 and the 40-update 1.065703117562344. Persistence is lower still: 0.4391345726797162. For glove angles, the final 0.37219431399606173 improves on 0.8750431745338253 and 1.1992171367822229, while persistence reaches 0.024541819466343505. Better than an earlier model is therefore not the same as better than the stated simple alternative. Both full-TRAIN streams lose that comparison; the private checkpoint remains unpromoted.

Separating the offsets explains why a favourable slice would mislead. RMS beats persistence at ten rows and loses at one. Glove loses at both. The first 4-window RMS score improves, but its glove score becomes worse than both the initial and 40-update states. Showing only that prefix or the favourable RMS offset would change the question after seeing the answers. The full comparison keeps them together, with the original 40-update failure intact.

The continuation changed several things together: batches grew from 4 to 32, there were more updates, and optimization revisited the full ordered catalog. We cannot assign the change to coverage alone. Five passes do not create five datasets, and one acquisition does not test whether a wider variety of people or tasks improves prediction. No DEV or TEST result is supplied. The accepted software and recorded arithmetic checks establish the bounded calculation, not independent generalization or a validated body model.

The broader programme still needs four different sources: neural observations from the relevant brain, spinal, peripheral and visceral systems; muscle measurements; measured movement; and physics-generated examples. Touch, proprioception and other task context need their own evidence. This diagnostic records none of those missing neural or sensory pathways and establishes no cross-modal mapping, learned force controller, biological teaching or scaling law. A model must first beat the relevant simple predictor for the question it actually answers. Testing whether a trainee retains a skill after trainer commands end remains a further experiment, with its own matched information, practice, assistance and movement-bridge controls.
